\documentclass[11pt]{article}
\usepackage[margin=1in]{geometry}
\usepackage{amsmath,amssymb}
\usepackage{booktabs}
\usepackage{enumitem}
\usepackage{caption}

\newcommand{\LE}{\mathrm{LE}}
\newcommand{\RE}{\mathrm{RE}}

\title{\textbf{Motivating Data Patterns} \\ Female Lifetime Earnings Inequality}
\author{}
\date{\today}

\begin{document}
\maketitle

\noindent This document summarizes key empirical patterns from the GKOS women's moments (SSA administrative data) that motivate and discipline the model in the accompanying specification. The data contain two tables: (i) cross-sectional moments of 1-year arc percent earnings changes by recent earnings (RE) group and age group (sheet \texttt{L1\_arc\_age\_re}), and (ii) mean earnings levels and growth rates by lifetime earnings (LE) percentile (sheet \texttt{incgrowth}). We organize the patterns separately for RE-based and LE-based moments.

%===============================================================================
\section{Recent Earnings (RE) Patterns}
%===============================================================================

The RE-based moments capture \emph{short-run earnings dynamics}---how annual earnings changes vary across the earnings distribution and over the lifecycle. RE groups are constructed from a 5-year trailing average of earnings, ranked within age group. The key statistic is the 1-year arc percent change, $\Delta^{\text{arc}} = 2(Y_{t+1} - Y_t)/(Y_{t+1} + Y_t)$, which is symmetric around zero and bounded by $[-2, 2]$.

\subsection{Pattern RE-1: Large negative left tail, concentrated at low RE}

\begin{table}[h]\centering
\caption{P10 of 1-year arc percent earnings change}
\begin{tabular}{lrrr}
\toprule
RE Group & Ages 25--34 & Ages 35--44 & Ages 45--54 \\
\midrule
Bottom 10 & $-1.59$ & $-1.42$ & $-1.42$ \\
10--30    & $-1.16$ & $-0.97$ & $-0.90$ \\
30--50    & $-0.87$ & $-0.64$ & $-0.56$ \\
50--70    & $-0.69$ & $-0.46$ & $-0.39$ \\
70--90    & $-0.60$ & $-0.37$ & $-0.31$ \\
Top 10    & $-0.62$ & $-0.48$ & $-0.41$ \\
\bottomrule
\end{tabular}
\end{table}

\noindent\textbf{Pattern.} The 10th percentile of earnings changes is deeply negative for low-RE women---a P10 of $-1.59$ means that 10\% of women in the bottom RE group experience near-total earnings collapse (arc percent change near $-2$ corresponds to going from positive to near-zero earnings). This left tail shrinks monotonically as RE rises, leveling off around the 70th--90th RE percentile.

\medskip\noindent\textbf{Model connection.} The model generates this through two channels that disproportionately affect low-RE women:
\begin{enumerate}[nosep]
  \item \textbf{Birth-induced separation:} Low-RE women are at low-productivity firms where $\pi_{\text{paid}}(p)$ and $\pi_{\text{unpaid}}(p)$ are small, so births more often trigger full separation and near-zero earnings. The within-year timing draw $\tau$ generates a \emph{continuous} left tail rather than discrete mass points.
  \item \textbf{Exogenous separation:} $\delta(p)$ is decreasing in $p$, so low-$p$ workers face higher baseline separation risk, amplifying the negative tail.
\end{enumerate}

\subsection{Pattern RE-2: Negative skewness everywhere, worsening with age at low RE}

\begin{table}[h]\centering
\caption{Kelley skewness of 1-year arc percent earnings change}
\begin{tabular}{lrrr}
\toprule
RE Group & Ages 25--34 & Ages 35--44 & Ages 45--54 \\
\midrule
Bottom 10 & $-0.18$ & $-0.22$ & $-0.29$ \\
10--30    & $-0.20$ & $-0.24$ & $-0.33$ \\
30--50    & $-0.25$ & $-0.26$ & $-0.35$ \\
50--70    & $-0.32$ & $-0.27$ & $-0.32$ \\
70--90    & $-0.36$ & $-0.26$ & $-0.29$ \\
Top 10    & $-0.35$ & $-0.26$ & $-0.29$ \\
\bottomrule
\end{tabular}
\end{table}

\noindent\textbf{Pattern.} Earnings changes are negatively skewed for \emph{all} women---the distribution has a longer left tail than right tail. Two gradients emerge:
\begin{itemize}[nosep]
  \item \emph{Age gradient at low RE:} Skewness becomes more negative with age for the bottom half of the RE distribution (e.g., $-0.18 \to -0.22 \to -0.29$ for the bottom 10).
  \item \emph{RE gradient at young ages:} Among young women (25--34), skewness is \emph{more} negative at the top of the distribution ($-0.36$ for RE 70--90) than the bottom ($-0.18$).
\end{itemize}

\medskip\noindent\textbf{Model connection.} The age gradient at low RE reflects the declining role of human capital growth as a positive shock: $\psi(\alpha, t)$ decreases with age, so the right tail of earnings changes shrinks, making the distribution more left-skewed. For young high-RE women, the strongly negative skewness reflects that they have ``more to lose''---they sit at high-$p$ firms with high $h$, but face the same birth-event risks. A birth-induced separation or amenity switch generates a large negative change relative to their high baseline. The job ladder (acceptance rule $p' > p$) limits upside moves for those already near the top.

\subsection{Pattern RE-3: Dispersion declines with RE and age}

\begin{table}[h]\centering
\caption{P90--P10 spread of 1-year arc percent earnings change}
\begin{tabular}{lrrr}
\toprule
RE Group & Ages 25--34 & Ages 35--44 & Ages 45--54 \\
\midrule
Bottom 10 & $2.74$ & $2.36$ & $2.20$ \\
10--30    & $1.94$ & $1.57$ & $1.35$ \\
30--50    & $1.38$ & $1.01$ & $0.83$ \\
50--70    & $1.04$ & $0.70$ & $0.58$ \\
70--90    & $0.87$ & $0.58$ & $0.46$ \\
Top 10    & $0.91$ & $0.74$ & $0.62$ \\
\bottomrule
\end{tabular}
\end{table}

\noindent\textbf{Pattern.} Earnings volatility is dramatically higher for low-RE women: the P90--P10 spread for the bottom 10 is roughly 3$\times$ that of the 70--90 group at every age. Both RE and age reduce dispersion, except that the top 10 shows slightly \emph{higher} dispersion than the 70--90 group (a ``U-shape'' at the top).

\medskip\noindent\textbf{Model connection.}
\begin{itemize}[nosep]
  \item \emph{RE gradient:} Low-RE women are at low-$p$ firms with high $\delta(p)$, high birth-separation risk, and the greatest scope for upward job-ladder moves ($p' > p$ is easily satisfied). This generates both large drops and large gains.
  \item \emph{Age gradient:} Human capital growth $\psi(\alpha, t)$ declines with age, reducing upside potential. Workers are further up the job ladder, reducing both separation risk (lower $\delta(p)$) and upward mobility (fewer firms above $p$).
  \item \emph{Top-10 uptick:} Partially reflects that the top RE group includes women with high $\alpha$ who may have birth events generate large percentage changes even at high earnings levels.
\end{itemize}

\subsection{Pattern RE-4: Young women face the most volatile earnings}

Across every RE group, the youngest age group (25--34) exhibits the highest P90--P10 spread and the deepest P10 values. For example, the P90--P10 spread for the 30--50 RE group drops from 1.38 at ages 25--34 to 0.83 at ages 45--54---a 40\% reduction.

\medskip\noindent\textbf{Model connection.} Ages 25--34 are the prime fertility years. The model predicts concentrated earnings disruptions in this window through birth events (both exogenous shocks and endogenous fertility choices), combined with still-rapid human capital accumulation that generates large upside moves. The three-outcome birth structure (paid leave, unpaid leave, separation) with within-year timing produces a wide range of outcomes for young workers that compresses as fertility tapers off after age 40.

%===============================================================================
\section{Lifetime Earnings (LE) Patterns}
%===============================================================================

The LE-based moments capture \emph{long-run earnings trajectories}---how the level and growth of earnings vary across the lifetime earnings distribution. LE is defined as $\sum_{t=25}^{55} w_t$ and women are ranked into percentiles by this cumulative measure.

\subsection{Pattern LE-1: Massive fan-out of earnings over the lifecycle}

\begin{table}[h]\centering
\caption{Mean earnings (\$2010) by LE group and age}
\begin{tabular}{lrrrr}
\toprule
LE Group & Age 25 & Age 35 & Age 45 & Age 55 \\
\midrule
Bottom 10 & \$5,683  & \$4,564  & \$6,912   & \$5,807 \\
10--30    & \$9,060  & \$8,968  & \$13,366  & \$13,496 \\
30--50    & \$12,248 & \$15,264 & \$22,006  & \$23,176 \\
50--70    & \$15,834 & \$23,112 & \$33,039  & \$34,911 \\
70--90    & \$20,300 & \$34,690 & \$50,468  & \$54,513 \\
Top 10    & \$26,176 & \$59,632 & \$102,964 & \$116,391 \\
\bottomrule
\end{tabular}
\end{table}

\noindent\textbf{Pattern.} The top-10/bottom-10 earnings ratio expands from 4.6$\times$ at age 25 to 20.0$\times$ at age 55---a 4.4-fold expansion of inequality over the working life. This fan-out is the central fact the model must explain for women.

\medskip\noindent\textbf{Model connection.} Three forces drive the fan-out:
\begin{enumerate}[nosep]
  \item \textbf{Job ladder:} High-$\alpha$ women accumulate human capital faster ($\psi_\alpha > 0$) and climb to higher-$p$ firms. Since wages are $w = p \cdot h - c(a)$, the multiplicative structure amplifies initial differences.
  \item \textbf{Differential fertility costs:} Low-LE women are more likely to experience birth-induced separations (low $\pi_{\text{paid}}$, low $\pi_{\text{unpaid}}$), which destroy firm matches and depreciate $h$. High-LE women at high-$p$ firms retain their jobs through paid leave.
  \item \textbf{Amenity sorting:} Low-LE women who become mothers at low-$p$ firms see the amenity cost $c(H)$ take a larger bite of their already-low wages, compounding the earnings gap.
\end{enumerate}

\subsection{Pattern LE-2: Earnings \emph{decline} for the bottom of the LE distribution}

\begin{table}[h]\centering
\caption{Log earnings growth by LE group and age window}
\begin{tabular}{lrrr}
\toprule
LE Group & Growth 25--35 & Growth 35--45 & Growth 45--55 \\
\midrule
Bottom 10 & $-0.23$ & $0.42$ & $-0.19$ \\
10--30    & $-0.02$ & $0.40$ & $0.00$ \\
30--50    & $0.21$  & $0.37$ & $0.05$ \\
50--70    & $0.37$  & $0.36$ & $0.06$ \\
70--90    & $0.53$  & $0.37$ & $0.08$ \\
Top 10    & $0.79$  & $0.49$ & $0.12$ \\
\bottomrule
\end{tabular}
\end{table}

\noindent\textbf{Pattern.} The bottom-10 LE group experiences \emph{negative} earnings growth from 25 to 35 ($-0.23$) and again from 45 to 55 ($-0.19$). This is striking: the decade when most women experience peak wage growth is a period of earnings \emph{decline} for the lowest-LE women. All groups share a mid-career hump (positive growth 35--45), but growth is monotonically increasing in LE at every age window.

\medskip\noindent\textbf{Model connection.}
\begin{itemize}[nosep]
  \item \emph{Decline at 25--35 for low LE:} This is the prime fertility window. Low-$\alpha$ women face high exogenous birth probability $\lambda_{\text{exo}}(\alpha)$ (decreasing in $\alpha$), combined with low-$p$ firms that offer neither paid nor unpaid leave. Birth events trigger separation, $h$ depreciation, and a reset to unemployment. Pre-25 births (from the early phase) compound the problem by giving these women lower $h_{25}$ and existing childcare needs.
  \item \emph{Recovery at 35--45:} The universal positive growth in this window reflects that most women have completed their fertility. The job ladder resumes its role: re-employed women climb from low $p$ back up. As children age ($k$ moves from 0--2 to 3--5 to 6+), amenity preferences shift back ($g(k) \to 0$), and women switch to $a = L$, recovering the wage cut $c(a)$.
  \item \emph{Late-career decline for low LE:} $\psi(\alpha, t)$ declines with age. For low-$\alpha$ women, human capital growth turns negligible or negative, while exogenous separation $\delta(p)$ at low-$p$ firms continues. This produces the ``last-decade fade'' in the bottom of the distribution.
\end{itemize}

\subsection{Pattern LE-3: Extreme earnings growth at the top}

The top-10 LE group experiences cumulative log earnings growth of 1.40 from age 25 to 55---a roughly 4-fold increase in earnings. Growth from 25 to 35 alone is 0.79, nearly four times the growth of the 10--30 group ($-0.02$) over the same window.

\medskip\noindent\textbf{Model connection.} High-$\alpha$ women benefit from the full compound effect of the job ladder:
\begin{itemize}[nosep]
  \item Rapid human capital accumulation ($\psi_\alpha \cdot \alpha$ large).
  \item Steady job-ladder climbing with low separation risk ($\delta(p)$ small at high $p$).
  \item When births occur, paid leave ($\pi_{\text{paid}}(p)$ high) preserves the firm match and $h$. The earnings dip is small and temporary.
  \item Amenity costs are affordable: $c(H)$ is a small fraction of $p \cdot h$ for high-$p$, high-$h$ women.
\end{itemize}

\subsection{Pattern LE-4: The 35--45 ``catch-up'' is universal but insufficient}

All six LE groups show positive growth in the 35--45 window, with remarkably similar rates (0.37--0.49). This uniformity contrasts sharply with the wide dispersion in the 25--35 window ($-0.23$ to $0.79$).

\medskip\noindent\textbf{Model connection.} By 35--45, most women have exited the fertile window (fertility declines steeply after 40). The ``catch-up'' growth at the bottom reflects:
\begin{itemize}[nosep]
  \item Re-entry to employment after birth-induced unemployment spells.
  \item Resumption of job-ladder climbing ($p' > p$ acceptance).
  \item Unwinding of amenity sorting as $g(k) \to 0$.
\end{itemize}
However, the catch-up is \emph{insufficient} to close the gap opened during 25--35, because: (i) $h$ was depreciated during unemployment, creating a permanent level shift; (ii) firm match quality was reset to a lower $p$, requiring years to rebuild; (iii) the job ladder is a slow process---only $\lambda_1$ fraction of workers receive offers each period.

\subsection{Pattern LE-5: Late-career deceleration is steeper at the bottom}

Growth from 45 to 55 is near zero or negative for the bottom half of the LE distribution, while the top decile still grows at 0.12. The bottom-10 group's $-0.19$ late-career decline means these women end their career earning \emph{less} than at age 45.

\medskip\noindent\textbf{Model connection.} The interaction of declining $\psi(\alpha, t)$ with the job ladder explains this: low-$\alpha$ women at low-$p$ firms still face meaningful separation risk ($\delta(p)$ high), but now have little human capital growth to compensate. High-$\alpha$ women at high-$p$ firms have negligible separation risk and still accumulate some $h$, maintaining modest positive growth even late in career.

%===============================================================================
\section{Cross-Cutting Patterns}
%===============================================================================

\subsection{Pattern X-1: Fertility as the key driver of divergence}

The sharpest divergence across LE groups occurs precisely during the fertility window (ages 25--35). The RE data confirm this: the youngest age group shows the highest dispersion and deepest left tail. Together, these patterns suggest that \emph{how women interact with the labor market around childbirth} is the primary source of lifetime earnings inequality---not initial differences in ability, which are modest at age 25 (4.6$\times$ top/bottom ratio).

\subsection{Pattern X-2: Institutional gradient in downside risk}

Across both RE and LE tables, the gradient in downside risk maps naturally onto the model's three-outcome birth structure:
\begin{itemize}[nosep]
  \item \textbf{High-RE/LE women} $\to$ high-$p$ firms $\to$ paid leave likely $\to$ job preserved, partial wage replacement $\to$ small earnings dip.
  \item \textbf{Mid-RE/LE women} $\to$ mid-$p$ firms $\to$ unpaid leave possible $\to$ job preserved but zero income during leave $\to$ moderate dip.
  \item \textbf{Low-RE/LE women} $\to$ low-$p$ firms $\to$ separation likely $\to$ loss of firm match, $h$ depreciation $\to$ near-total earnings collapse.
\end{itemize}
This maps directly to the institutional structure: paid employer leave at large firms, FMLA-eligible unpaid leave at medium firms, and no protection at small firms.

\subsection{Pattern X-3: Partial but incomplete recovery}

Both the RE and LE data show that negative shocks are partially---but not fully---reversed. The LE data show a universal mid-career rebound (Pattern LE-4), while the RE data show that P10 values are less extreme at older ages (Pattern RE-1). However, the LE fan-out (Pattern LE-1) demonstrates that the gap opened during the fertility years persists. The model accounts for this through the permanent destruction of firm match quality upon separation: even after re-employment, a woman starts at a lower $p$ and must re-climb the ladder, a process that takes years and may never reach the counterfactual path.

\end{document}
